% ECONOMICS_PROBLEM: {"id": "D31-162", "title": "Drivers of extreme wealth concentration", "jel_code": "D31", "source_id": "OP-0362"}
% !TeX program = xelatex
\documentclass[11pt,a4paper]{article}
\usepackage[margin=23mm]{geometry}
\usepackage{fontspec}
\setmainfont{Latin Modern Roman}
\usepackage{amsmath,amssymb,mathtools}
\usepackage{xcolor,enumitem,booktabs,longtable,array,xurl,hyperref}
\definecolor{muted}{HTML}{53636C}
\hypersetup{colorlinks=true,urlcolor=blue,linkcolor=blue}
\setlength{\parindent}{0pt}
\setlength{\parskip}{5pt}
\newcommand{\R}{\mathbb R}
\newcommand{\E}{\mathbb E}
\newcommand{\N}{\mathbb N}
\newcommand{\ind}{\mathbf 1}
\newcommand{\Var}{\operatorname{Var}}
\newcommand{\Cov}{\operatorname{Cov}}
\newcommand{\supp}{\operatorname{supp}}
\DeclareMathOperator*{\argmin}{arg\,min}
\DeclareMathOperator*{\argmax}{arg\,max}
\newcommand{\entrymeta}[3]{{\sffamily\small\color{muted}\textbf{\texttt{#1}}\quad|\quad #2\par #3\par}}
\newcommand{\Problemref}[1]{\texttt{#1}}
\newcommand{\JELref}[2]{\texttt{#2} (JEL \texttt{#1})}
\begin{document}
\section*{Drivers of extreme wealth concentration}
\textbf{Problem ID:} \texttt{D31-162}\quad \textbf{JEL:} \texttt{D31}\par
% BEGIN ECONOMICS STATEMENT
\label{problem:OP-0362}
\label{audited:ECON-028}
{\sffamily\small\color{muted}Original subject group: Inequality and economic mobility\par}
\label{definitions:problem:30007644}
\textbf{Audit classification:} Published research agenda. Working text and published DOI metadata checked. Joint quantitative importance of saving mechanisms is future work; the structural attribution rule is editorial.
\subsection*{Definitions and precise research target}
\textbf{Complete editorial problem.} Consider a specified economy of households whose wealth evolves according to \(W_{i,t+1}=(1+r_{it})W_{it}+Y_{it}-C_{it}+B_{it}\), where \(r_{it}\) is the realized portfolio return, \(Y_{it}\) labor and entrepreneurial income not already included in portfolio returns, \(C_{it}\) consumption, and \(B_{it}\) net transfers including inheritances. Let \(\mathcal M\) be a declared class of models governing saving, portfolio choice, entrepreneurship, and intergenerational transfers. For a top population share \(p\in(0,1)\), assuming positive finite total wealth under every compared model, let \(S_{p,t}(M)\) be the model-implied share of total wealth owned by that group. Identify which models jointly reproduce concentration, individual wealth transitions, returns, portfolios, and inheritance histories, rather than merely matching \(S_{p,t}\). For each mechanism \(k\), define a counterfactual \(M^{-k}\) that removes that mechanism according to a stated intervention while recomputing equilibrium. Estimate \(\Delta_{k,t}=S_{p,t}(M)-S_{p,t}(M^{-k})\). Since mechanisms interact, report the counterfactual ordering or a specified averaging rule and avoid interpreting contributions as an invariant additive identity. The source-motivated quantitative task is to distinguish relative importance and validate predictions outside fitted moments. Assumptions about unobserved skill, bequest motives, and endogenous returns must be exposed and tested where possible.

\textbf{Definitions and admissible scope.} The wealth recursion uses a net financial return only when that aggregate portfolio return is well defined; with leverage, specify asset returns and debt interest separately. \(Y_{it}\) excludes profits and valuation changes already included in returns, and \(B_{it}\) is net of transfers paid to others. Define an admissible structural model by household preferences, earnings and return processes, borrowing constraints, birth/death and inheritance rules, prices, and equilibrium selection. Wealth top shares use quantile integrals, finite absolute wealth, and positive mean wealth. Removing mechanism \(k\) specifies the replacement primitive and the fiscal/resource closure; it is not merely setting an estimated coefficient to zero. A contribution depends on that intervention, and counterfactual interactions can prevent additive channel shares. Specify the distribution of baseline households and a fixed horizon \(t\). For every admissible model define \(S_{p,t}(M)=\int_{1-p}^1F^{-1}_{W_t,M}(u)\,du/E_M[W_t]\), with \(E_M|W_t|<\infty\) and \(E_MW_t>0\). The stated scalar-return recursion applies only when \(r_{it}W_{it}\) is a meaningful net portfolio payoff; a zero or leveraged net position requires separate asset and debt accounts. The class \(\mathcal M_0\) fixes the structural primitives, allowed parameter bounds, exogeneity restrictions, and selection rules. Restrict to its members matching the observed joint law. Each intervention \(M^{-k}\) is defined on those same baseline resources with an explicit replacement primitive and fiscal closure; return the joint identified set of \((\Delta_{k,t})_k\), not an invariant causal channel identity.
\subsection*{Variants and assumption changes}
\begin{enumerate}
\item Returns and endogenous portfolios (source-motivated): compare a common risk-adjusted return primitive with heterogeneous returns generated by knowledge and portfolio choice. Require the models to fit individual transitions as well as top shares.
\item Joint channel attribution (editorial): for \(K<\infty\) mechanisms define a model for every subset and average each mechanism's marginal contribution over all orderings. This Shapley-style accounting depends on the specified subset interventions and does not identify mechanisms without causal restrictions.
\end{enumerate}
\subsection*{References and source audit}
\textbf{Support and access.} Working text and published DOI metadata checked. Joint quantitative importance of saving mechanisms is future work; the structural attribution rule is editorial. \textbf{Locator:} CEPR DP 11746, Sections 8.1 and 10, pp. 26--27 and 32--33.
\begin{enumerate}\raggedright
\item\hypertarget{definitions:source:30007644:1}{} Mariacristina De Nardi; Giulio Fella. 2017. \emph{Saving and Wealth Inequality}. CEPR DP 11746; introduction, printed pp. 2–4; section 8.1, printed pp. 26–27.
\url{https://gabriel-zucman.eu/files/teaching/DeNardiFella17.pdf}.
DOI: \url{https://doi.org/10.1016/j.red.2017.06.002}.
\end{enumerate}

\subsection*{Integrated refinement: ECON-028}
{\sffamily\small\color{muted}Top-wealth dynamics with heterogeneous returns\par}
This refinement adopts the additional source conventions
(page~\pageref{audited:conventions}), subject to its local restrictions.
\textbf{Return opportunities and transition paths of top wealth shares.}
For a wealth law $F_t$ with finite absolute first moment and positive mean, use
\[
S_p(F_t)=\frac{\int_{1-p}^{1}F_t^{-1}(u)\,du}
                   {\int_0^1F_t^{-1}(u)\,du},\qquad 0<p<1.
\]
This assigns a cutoff atom fractionally and selects exactly the top $p$ share.
Specify taxes, mortality, entry, debt and inheritance in the wealth transition.
Identify a declared return-opportunity intervention's effect on the entire
$S_p(F_t)$ path, allowing saving, earnings, prices and demographics to adjust.
Equal mean returns, equal realized returns and equal access are different
interventions. Preserve the distinction between gross assets, net and liquid
wealth; tail indices, rank mobility and convergence speeds are separate targets.
A stationary tail or an accounting wealth recursion does not identify these
causal transition effects.
\subsection*{Supplied source audit for ECON-028}
Local [S1], [S2], etc. in this audit and the references immediately below
belong to ECON-028, independently of the original entry's reference numbering.
[S1], PDF p. 8, supports heterogeneous returns and wealth dynamics. The intervention and functional are editorial; the quantile-atom error is corrected.
\subsection*{References for integrated refinement ECON-028}
{\footnotesize\raggedright
\par\noindent\textbf{[S1]} Benjamin Moll (n.d.). \emph{Lecture 11: Good Research Topics and Open Questions}. Distributional Macroeconomics lecture slides; publicly archived at a 2019 URL.\par \textit{Verified locator / source role:} PDF p. 8: wealth inequality, return heterogeneity, and transition dynamics.\par \url{https://benjaminmoll.com/wp-content/uploads/2019/07/Lecture11_2149.pdf}\par\smallskip
}

\subsection*{Common conventions: Source mathematical conventions}

These conventions apply to each entry unless explicitly replaced.

\textbf{Models and identification.} A primitive model \(m\) specifies all probability laws, feasible choices, information, equilibrium and measurement rules used by its target. The admissible class \(\mathcal M\) is fixed independently of outcomes. If \(P_O(m)\) is its observable probability law and \(\tau(m)\) the target, its identified set at observable law \(P\) is
\[
 \mathcal I_\tau(P)=\{\tau(m):m\in\mathcal M,\ P_O(m)=P\}.
\]
Point identification means that this set is a singleton. An empty set signals incompatibility of data and assumptions. Coordinate bounds need not describe the joint set. Bounds are sharp only if every retained target is attainable or arbitrarily approachable by compatible models. Finite bounds require explicit restrictions on unobserved quantities; unrestricted confounding, hidden wealth, extrapolation or arbitrary equilibrium selection can yield uninformative sets.

\textbf{Causal interventions.} A potential outcome \(Y_i(p)\) is the outcome under a complete policy regime \(p\), including timing, financing, information and market spillovers. The target population and units are fixed before treatment. With interference, use the complete assignment vector or a justified exposure mapping; individual no-interference notation alone cannot identify an economy-wide intervention. A difference of mean potential outcomes does not require the cross-world joint distribution. Conditioning on post-treatment migration, survival, enrollment or employment changes the estimand and needs separate justification.

\textbf{Statistical statements.} An estimator, confidence set, forecast or test specifies a sampling law, model class and loss/coverage criterion. Uniform \((1-\alpha)\)-coverage means \(\inf_{m\in\mathcal M}\Pr_m\{\tau(m)\in C_n\}\geq1-\alpha\), or an explicitly asymptotic version. The whole parameter space trivially has coverage, so informativeness requires a stated width, risk or power objective. A forecasting improvement is relative to a named benchmark, information set and evaluation population.

\textbf{Equilibrium and optimization.} An equilibrium comprises feasible plans, optimality, beliefs where needed, budget constraints and all market/resource clearing. The primitives or a stated benchmark define each correspondence \(\mathcal E(p;m)\). Nonemptiness is assumed or proved, never inferred just from compact policy sets. A selected-equilibrium target uses a declared selection rule; a robust target quantifies over all permitted equilibria. Use suprema or infima unless attainment is established. An \(\varepsilon\)-optimal policy is feasible and within \(\varepsilon\) of the finite optimum value. Report an empty feasible set separately.

\textbf{Welfare and accounting.} Normative utility, interpersonal normalization, social weights, numeraire, discounting and the use of public receipts are primitives. Behavioral utility need not coincide with the welfare criterion. Household welfare includes ownership income and rents once; adding profits again to compensating variation generally double-counts them. A monetary equivalent is defined by an explicit transfer and price convention, with monotonicity and range conditions. Average effects, mechanism contributions, transfers and welfare losses are different targets. Infinite sums require convergence and dynamic feasible plans require terminal or no-Ponzi conditions.

\textbf{Source versions.} A working-paper date, revision date and journal publication year are distinguished. A DOI for a journal version is not automatically the DOI of an earlier manuscript. Locators refer to the stated version and printed pages where specified. A metadata-only check does not verify a claim about a full-text conclusion. Every entry's source subsection narrows the provenance claim accordingly.

\subsection*{Common conventions: Additional-source status and mathematical conventions}

\label{audited:conventions}

Of the 100 supplied ECON entries, 65 distinct problems appear here; 32 close
matches refine earlier entries, and three duplicate questions map to existing
entries. Appendix~\ref{app:audited-crosswalk} lists every destination.
The attachment's P/R/S/U distinctions and source audits are retained. Its
precise mathematical formalizations are editorial unless the individual audit
says otherwise. Candidate subsidy targets ECON-004--006 retain their U flags;
the resolved approval-core question ECON-007 updates OP-0202 above.
The following supplied conventions govern both this part and the attributed
ECON refinements elsewhere in the catalogue, subject to local restrictions.

\textbf{Parameterized questions.} A problem family lists inputs that must be fixed before an instance is evaluated: population, horizon, policies, information, data law, model class and welfare weights. These are required choices, not quantities the citations are assumed to specify. Undefined applications should not be treated as identified estimands.

\textbf{Indices and integrability.} Sets of agents, firms and regions are finite unless a population measure is explicitly used. \(i,j\) index units, \(r\) regions, \(t\) dates and \(h\) horizons. \(\mathbb E\) and \(\Pr\) refer to the declared population and law; \(\mathbf1\{A\}\) is an indicator. Every displayed expectation and discounted sum needs existence in the stated domain. Infinite horizons use a discount in \((0,1)\) and a convergence condition; finite horizons may use discount one.

\textbf{Optimization and existence.} A displayed \(\max\), \(\min\), or \(\arg\max\) is conditional on attainment. For finite feasible menus this is immediate; general spaces need continuity and compactness/coercivity or another existence argument. Without attainment, the target is the supremum/infimum and arbitrarily near-optimal feasible choices. \(\inf\varnothing=+\infty\). Empty feasibility and an unidentified target are different issues.

\textbf{Potential outcomes.} \(Y_i(z)\) is the outcome under a specified intervention, not the observed conditional mean among people choosing \(z\). In binary comparisons \(\Delta Y_i=Y_i(1)-Y_i(0)\). Specify assignment, treatment versions, compliance, information and observation horizon. Interference is allowed under a full policy vector; compressed exposure notation needs a justified mapping. No interference, consistency and overlap are substantive assumptions, not automatic consequences of notation.

\textbf{Populations and observation.} Fix baseline eligibility and population weights. Conditioning on post-policy employment, survival, relocation or program uptake can change the population being compared. Death is not ordinary loss to follow-up; outcomes truncated by death need a composite convention, joint survival target, or explicitly defined survivor stratum. Fertility-changing interventions need a population functional rather than an unexplained fixed descendant cohort. Measurement and missingness models must be supplied.

\textbf{Identification.} A structural model \(m\in\mathcal M\) implies observable law \(P_m\) and target \(\theta(m)\). Identification means
\[
 P_{m_1}=P_{m_2}\ \Longrightarrow\ \theta(m_1)=\theta(m_2)
 \quad\text{for all }m_1,m_2\in\mathcal M .
\]
Without identification, the target is
\[
 \Theta(P;\mathcal M)=\{\theta(m):m\in\mathcal M,\ P_m=P\}.
\]
Writing \(\theta(P)\) before establishing identification can hide incompatible latent models. Confidence sets additionally need a sampling law, confidence level, asymptotic sequence and uniformity class. Point identification, coverage and informative precision are separate properties.

\textbf{Mediation.} Total effects of randomized assignment are distinct from cross-world natural indirect effects. Belief/effort-channel effects require a meaningful mediator intervention and additional measurement, exclusion or mediation assumptions. Randomization of the initial policy alone does not supply those assumptions. Report bounds or interventional alternatives when the stronger target is unsupported.

\textbf{Equilibrium.} A counterfactual \(W(z)\), \(p(z)\), or \(\mu_t^z\) requires individual optimization, feasibility, government/resource budgets, market clearing and state-transition laws. State equilibrium existence and selection, or report ranges over compatible equilibria. Initial-state integration includes subsequent endogenous transitions; current-state integration must use the policy-dependent distribution. Reduced-form invariance after policy is not assumed.

\textbf{Welfare accounts.} Scores, utility, health, dollars and ecological quantities require explicit conversion before addition. Fees, taxes, subsidies, wages and extortion payments are transfers between parties; distributional weights can make them welfare relevant, but their face value is not automatically a resource loss. State ownership, geographic boundaries, foreign-income weights and fiscal finance. Do not add profits to household welfare already including dividends, or count land capitalization and the underlying benefit twice. A benefit-cost expression must distinguish gross benefits from net welfare and subtract every real cost exactly once.

\textbf{Derivatives and risk.} Log derivatives require positive arguments; local responses require admissible perturbations and specified equilibrium branches. Differentiation under expectations needs regularity/domination, and boundaries may allow only one-sided derivatives. Risk uses a specified law; ambiguity uses a declared nonempty law set on a common history space. Adaptive policies must be nonanticipative. A robust criterion is an editorial choice unless attributed explicitly.

\textbf{Scope overlaps.} Allocation, collective choice, contracts, household bargaining and coordinated policy overlap economics with optimization and game theory. They are retained because they appeared in the original catalogue; no claim of a strictly game-theory-free selection is made.
% END ECONOMICS STATEMENT
\end{document}
